\documentclass[12pt, letterpaper]{article}

%-------Packages---------
\usepackage{amssymb,amsfonts}
\usepackage{mathptmx}
\usepackage{amsthm}
\usepackage{graphicx}
\usepackage[all,arc]{xy}
\usepackage[colorlinks=true, allcolors=blue]{hyperref}
\usepackage{enumerate}
\usepackage{bbm}
\usepackage{dsfont}
\usepackage{mathrsfs}
\usepackage{tikz-cd}
\usepackage{setspace}
\usepackage{import}
\usepackage[utf8]{inputenc}
\usepackage{hyperref}
\usepackage{tikz}
\usepackage{float}
\usepackage{mdframed}
\usepackage{fancyhdr}
\usepackage[nointegrals]{wasysym}

\usepackage[left=1in,top=1in,right=1in,bottom=1in]{geometry}

\usepackage{mathtools}
\usepackage{setspace}
\usepackage{cleveref}
\usepackage{multicol}

%--------Theorem Environments--------
%theoremstyle{plain} --- default
\newmdtheoremenv{theorem}{Theorem}[subsection]
\newmdtheoremenv{corollary}[theorem]{Corollary}
\newtheorem{proposition}[theorem]{Proposition}
\newmdtheoremenv{lemma}[theorem]{Lemma}
\newtheorem{conj}[theorem]{Conjecture}
\newtheorem{quest}[theorem]{Question}

\theoremstyle{definition}
\newmdtheoremenv{definition}[theorem]{Definition}
\newmdtheoremenv{definitions}[theorem]{Definitions}
\newtheorem{example}[theorem]{Example}
\newtheorem{algorithm}[theorem]{Algorithm}
\newtheorem{rmk}[theorem]{Remark}
\newtheorem{exercise}[theorem]{Exercise}

%----------- my commands -------------
\newcommand{\C}{\mathbb{C}}
\newcommand{\R}{\mathbb{R}}
\newcommand{\Q}{\mathbb{Q}}
\newcommand{\Z}{\mathbb{Z}}
\newcommand{\N}{\mathbb{N}}
\renewcommand{\P}{\mathbb{P}}
\newcommand{\contr}{\Rightarrow \Leftarrow}
\newcommand{\HRule}[1]{\rule{\linewidth}{#1}}
\newcommand{\s}{\(\,\)}
\renewcommand{\t}[1]{\text{#1}}
\newcommand{\ang}[1]{\left\langle #1 \right\rangle}

% Meta data
\setstretch{1.2}

\begin{document}

\pagestyle{fancy}
\fancyhf{}
\fancyhead[LOE]{\slshape{\textbf{Vincent Ferrara}}}
\fancyhead[ROE]{\thepage}

\subsection*{Problem 1}
\begin{itemize}
    \item If $n$ is odd
    \begin{itemize}
        \item Consider $x^i$ for $0 < i < n$ s.t. $x$ is a rotation\\
        $x^jx^ix^{-j}$ for $0 \leq j < n$\\
        $=x^i$\\
        $yx^jx^iyx^{j}$ for $0 \leq j < n$\\
        $=yyx^{n-j-i}x^{j}=x^nx^{-i}x^{j-j}=x^{-i}$\\
        Thus $C_{x^{i}}=\{x^i,x^{-i}\}=C_{x^{-i}}$
        \item Consider $y$ s.t. $y$ is a reflection\\
        $x^jyx^{-j}$ for $0 \leq j < n$\\
        $=yx^{n-j}x^{-j}=yx^{-2j}$\\
        Since $n$ is odd we can solve $k \equiv -2j \pmod{n}$\\
        thus $C_y=\{yx^k \mid k \in \Z\}=\{yx^k \mid k \in \{0,1,2,\dots,n-1\}\}$
        
        Thus $|D_n|=\sum\limits_{[x^i]}|C_{x^i}|+|C_y|=|\{e\}|+\sum\limits_{[x^i]\setminus [e]}|C_{x^i}|+|C_y|=1+(n-1)+n=2n$
    \end{itemize}
    \item If $n$ is even
    \begin{itemize}
        \item It is the same for rotations
        \item Consider $y$ s.t. $y$ is a reflection\\
        $x^jyx^{-j}$ for $0 \leq j < n$\\
        $=yx^{n-j}x^{-j}=yx^{-2j}$\\
        thus $C_y=\{y,yx^2,yx^4,\dots , yx^{n/2 - 1}\}$
        \item Consider $yx$\\
        $x^jyxx^{-j}$ for $0 \leq j < n$\\
        $=yx^{n-j}x^{1-j}=yx^{1-2j}$\\
        thus $C_{yx}=\{yx,yx^3,yx^5,\dots , yx^{n-1}\}$

        Thus $|D_n|=\sum\limits_{[x^i]}|C_{x^i}|+|C_y|+|C_{yx}|$

        $=|\{e\}|+|\{x^{n/2}\}|+\sum\limits_{[x^i]\setminus [e]\cup[x^{n/2}]}|C_{x^i}|+|C_y|+|C_{yx}|=1+1+2(n/2-1)+2(n/2)=2n$
    \end{itemize}
\end{itemize}

\newpage
\subsection*{Problem 2}
\begin{itemize}

    \item Conjugacy class of \( 1 \):
    \[
    \{ 1 \}
    \]
    Since \( 1 \) is the identity element, it commutes with all other elements:
    \[
    h \cdot 1 \cdot h^{-1} = 1 \quad \text{for all} \quad h \in Q_8.
    \]

    \item Conjugacy class of \( -1 \):
    \[
    \{ -1 \}
    \]
    Similar to \( 1 \), the element \( -1 \) commutes with all other elements:
    \[
    h \cdot (-1) \cdot h^{-1} = -1 \quad \text{for all} \quad h \in Q_8.
    \]

    \item Conjugacy class of \( i \) and \( -i \):
    
    The conjugates of \( i \) by other elements of \( Q_8 \) are:
    \[
    1 \cdot i \cdot 1^{-1} = i, \quad (-1) \cdot i \cdot (-1)^{-1} = i,
    \]
    \[
    j \cdot i \cdot j^{-1} = -i, \quad k \cdot i \cdot k^{-1} = -i.
    \]
    Therefore, the conjugacy class of \( i \) is:
    \[
    \{ i, -i \}.
    \]
    
    Similarly, for \( -i \):
    \[
    j \cdot (-i) \cdot j^{-1} = i, \quad k \cdot (-i) \cdot k^{-1} = i,
    \]
    so \( i \) and \( -i \) share the same conjugacy class.

    \item Conjugacy class of \( j \) and \( -j \):
    
    The conjugates of \( j \) are computed as follows:
    \[
    1 \cdot j \cdot 1^{-1} = j, \quad (-1) \cdot j \cdot (-1)^{-1} = j,
    \]
    \[
    i \cdot j \cdot i^{-1} = -j, \quad k \cdot j \cdot k^{-1} = -j.
    \]
    Hence, the conjugacy class of \( j \) is:
    \[
    \{ j, -j \}.
    \]
    
    Similarly, for \( -j \):
    \[
    i \cdot (-j) \cdot i^{-1} = j, \quad k \cdot (-j) \cdot k^{-1} = j,
    \]
    so \( j \) and \( -j \) are in the same conjugacy class.

    \item Conjugacy class of \( k \) and \( -k \):
    
    The conjugates of \( k \) are:
    \[
    1 \cdot k \cdot 1^{-1} = k, \quad (-1) \cdot k \cdot (-1)^{-1} = k,
    \]
    \[
    i \cdot k \cdot i^{-1} = -k, \quad j \cdot k \cdot j^{-1} = -k.
    \]
    Therefore, the conjugacy class of \( k \) is:
    \[
    \{ k, -k \}.
    \]
    
    Similarly, for \( -k \):
    \[
    i \cdot (-k) \cdot i^{-1} = k, \quad j \cdot (-k) \cdot j^{-1} = k,
    \]
    so \( k \) and \( -k \) share the same conjugacy class.

\end{itemize}
\phantom{text}\\
Thus, the conjugacy classes of \( Q_8 \) are:
\[
\{ 1 \}, \quad \{ -1 \}, \quad \{ i, -i \}, \quad \{ j, -j \}, \quad \{ k, -k \}.
\]
The sizes of these conjugacy classes are:
\[
| \{ 1 \} | = 1, \quad | \{ -1 \} | = 1, \quad | \{ i, -i \} | = 2, \quad | \{ j, -j \} | = 2, \quad | \{ k, -k \} | = 2.
\]

Therefore, the class equation of \( Q_8 \) is:
\[
8 = 1 + 1 + 2 + 2 + 2.
\]

\newpage
\subsection*{Problem 3}
\begin{enumerate}[(a)]
    \item \begin{itemize}
        \item (i. $\implies$ iii.) Assume i.\\
        Let $\varphi:G \rightarrow K$ s.t. $\varphi(g)=\varphi(hk)=k$\\
        Well Defined:
        Let $g \in G$ by i. since $g=hk$ is unique this implies every input has one output.\\
        Kernel:
        Let $g \in$ ker$(\varphi)$\\
        $\varphi(g)=\varphi(hk)=e \implies k=e$ thus $g=h \in H$\\
        Let $h \in H$ thus $he \in HK$ thus $\varphi(h)=\varphi(he)=e$ thus $H=$ker$(\varphi)$\\
        Surjective: Let $k \in K$ consider: $ek \in G$\\
        $\varphi(ek)=k$ thus $\varphi$ is surjective.\\
        By the first isomorphism theorem this implies $K \cong G\diagup H$
        \item (iii. $\implies$ ii.) Assume iii. with the map $\varphi$\\
        Since $\varphi$ is injective ker$(\varphi)=\{e\} \implies$ ker$(\varphi) \subseteq K \cap H$
        Let $l \in K \cap H \implies l \in K$ and $l \in H$ thus $\varphi(l)=lH$ but since $l \in H \implies \varphi(l)=H$ thus $l \in$ ker$(\varphi)$ thus ker$(\varphi)=K \cap H=\{e\}$\\
        Since $\varphi$ is surjective this means for all $kH \in G \diagup H$ there exists a $k' \in K$ s.t. $\varphi(k')=kH$ and since $G \diagup H$ paritions $G$ this implies for all $g \in G \implies g \in jH$ for some $j \in K \implies g=jh$ for some $h \in H$\\
        thus $G=HK$
        \item (ii. $\implies$ i.) Assume ii.\\
        Let $g \in G$ s.t. $g=hk=ij$ for $h,i \in H$ and $kj \in K$\\
        thus $i^{-1}h=jk^{-1} \implies$ that is in $H \cap K$ thus $i^{-1}h=e \implies i=h$ and $jk^{-1}=e \implies j=k$ thus all $g\in G$ are a unique product of an $h \in H$ and $k \in K$ 
    \end{itemize}
    \item Let $\varphi : G \rightarrow (H \times K, \t{ with the given group multiplication})$ s.t. $\varphi(hk)=(h,k)$\\
    Homomorphism: Let $g_1,g_2 \in G$ thus $g_1=h_1k_1$ and $g_2=h_2k_2$ for $h_1,h_2 \in H$ and $k_1,k_2 \in K$\\
    $\varphi(g_1g_2)=\varphi(h_1k_1h_2k_2)=\varphi(h_1(k_1h_2k_1^{-1})k_1k_2)$ since $H$ is normal in $G$ this means $k_1h_2k_1^{-1} \in H$\\
    thus $\varphi(h_1(k_1h_2k_1^{-1})k_1k_2)=(h_1(k_1h_2k_1^{-1}), k_1k_2)=(h_1,k_1)(h_2,k_2)=\varphi(g_1)\varphi(g_2)$\\
    Injective: Let $g \in$ ker$(\varphi)$ thus $g=hk$ for $h \in H$ and $k \in K$\\
    $\varphi(g)=\varphi(hk)=(h,k)=(e,e)$ thus $h=e$ and $k=e$ thus $g=ee=e$ thus ker$(\varphi)=\{e\}$\\
    Surjective: Let $(h,k) \in (H \times K, \t{ with the given group multiplication})$ consider $hk=g \in G$\\
    $\varphi(g)=\varphi(hk)=(h,k)$\\
    Thus $G \cong (H \times K, \t{ with the given group multiplication})$
    \item Let $k \in K$ consider $\phi_k$\\
    Homomorphism: Let $h,j \in H$\\
    $\phi_k(hj)=khjk^{-1}=khk^{-1}kjk^{-1}=\phi_k(h)\phi_k(j)$\\
    Injective: Let $h \in$ ker$(\phi_k)$\\
    $\phi_k(h)=khk^{-1}=e \implies h=e$ thus ker$(\phi_k)=\{e\}$\\
    Surjective: Let $h \in H$ since $H$ is normal conisder $k^{-1}hk$\\\
    $\phi_k(k^{-1}hk)=kk^{-1}hkk^{-1}=h$\\
    Thus $\phi_k$ is an automorphism\\
    \phantom{text}\\
    Homomorphism: Let $k,j \in K$\\
    $\phi(kj)(h)=\phi_{kj}(h)=kjhj^{-1}k^{-1}=\phi_k\circ\phi_j(h)=\phi(k)\phi(j)(h)$\\
    Thus $\phi$ is a homomorphism\\
    \phantom{text}\\
    If $H,K$ are an internal direct product then $\phi(k)=id_H$ for all $k\in H$\\
    Thus the group law of $H \times K$ would be $(h_1,k_1)(h_2,k_2)=(h_1h_2,k_1k_2)$ thus proven in last homework this means $G$ is a direct product of $H,K$
\end{enumerate}

\newpage
\subsection*{Problem 4}
\begin{enumerate}[(a)]
    \item \begin{itemize}
        \item (Closure) Given by group law this is closed
        \item (Identity) Consider $(e,e),(h,k)\in H \rtimes_{\phi} K$\\
        $(e,e)(h,k)=(e\phi_e(h),ek)$ since $\phi$ is a homomorphism and $\phi(e)=\phi_e$ this means that $\phi(e)=\phi_e=id_H$
        thus $=(h,k)$\\
        $(h,k)(e,e)=(h\phi_k(e),ke)$ since $\phi_k$ is an automorphism it is a homomorphism\\
        thus $\phi_k(e)=e$ thus $=(he,k)=(h,k)$
        \item (Associativity) Let $(h_1,k_1),(h_2,k_2),(h_3,k_3) \in H \rtimes_\phi K$\\
        $(h_1,k_1)[(h_2,k_2)(h_3,k_3)]=(h_1,k_1)(h_2\phi_{k_2}(h_3),k_2k_3)=(h_1\phi_{k_1}(h_2\phi_{k_2}(h_3)),k_1k_2k_3)$\\
        $[(h_1,k_1)(h_2,k_2)](h_3,k_3)=(h_1\phi_{k_1}(h_2),k_1k_2)(h_3,k_3)=(h_1\phi_{k_1}(h_2)\phi_{k_1k_2}(h_3),k_1k_2k_3)\\
        =(h_1\phi_{k_1}(h_2\phi_{k_2}(h_3),k_1k_2k_3))$
        \item (Inverses) Let $(h,k) \in H \rtimes_\phi K$\\
        Consider $(\phi_{k^{-1}(h^{-1})},k^{-1})$\\
        $(h,k)(\phi_{k^{-1}}(h^{-1}),k^{-1})=(h\phi_{k}(\phi_{k^{-1}}(h^{-1})),kk^{-1})=(h\phi_{e}(h^{-1}),e)=(e,e)$\\
        $(\phi_{k^{-1}}(h^{-1}),k^{-1})(h,k)=(\phi_{k^{-1}}(h^{-1})\phi_{k^{-1}}(h),k^{-1}k)=(\phi_{k^{-1}}(e),e)=(e,e)$
    \end{itemize}
    \item The internal and external semidirect products provide two perspectives on how groups can combine through an action. 
    In the external semidirect product, we start with two independent groups \(H\) and \(K\) and define a homomorphism 
    \(\phi: K \to \text{Aut}(H)\) to describe how \(K\) acts on \(H\). 
    This action defines a group structure on the set \(H \times K\), where the multiplication is given by 
    \((h_1, k_1) \cdot (h_2, k_2) = (h_1 \cdot \varphi(k_1)(h_2), k_1 \cdot k_2)\). 
    The internal semidirect product, on the other hand, starts with a single group \(G\) that can be decomposed into two subgroups \(H \triangleleft G\) and \(K \leq G\). 
    Here, \(H\) is a normal subgroup, and \(K\) acts on \(H\) via conjugation: \(khk^{-1} \in H\) for \(k \in K\) and \(h \in H\). 
    Both constructions describe the same type of group structure, where the elements of one group act on another.
    \item \begin{itemize}
        \item (Identity) Let $\phi \in$ Hom$(K, \t{Aut}(H))$, and let $(\psi,\mu)=(id_H,id_K)$\\
        $((id_H,id_k)\phi)(k)=id_H \circ \phi(id_K^{-1}(k))\circ id_H^{-1}$\\
        $=id_H \circ \phi(k) \circ id_H = \phi(k)$
        \item Let $(\psi_1, \mu_1),(\psi_2,\mu_2) \in \t{Aut}(H) \times \t{Aut}(K)$ and let $\phi \in$ Hom$(K, \t{Aut}(H))$\\
        $[(\psi_1\circ\psi_2,\mu_1\circ\mu_2)\phi](k)=\psi_1\circ\psi_2\circ(\phi(\mu_2^{-1}\circ\mu_1^{-1}(k)))\circ\psi_2^{-1}\circ\psi_1^{-1}$\\
        $((\psi_1, \mu_1)[(\psi_2,\mu_2)\phi])(k)=(\psi_1, \mu_1)(\psi_2 \circ(\phi(\mu_2^{-1}(k)))\circ \psi_2^{-1})=\\
        \psi_1 \circ \psi_2 \circ(\phi(\mu_2^{-1}\circ\mu_1^{-1}(k)))\circ \psi_2^{-1} \circ \psi_1^{-1}$\\
        Thus this is an action
    \end{itemize}
    Assume $\phi, \phi'$ are in the same orbit for this action this means\\
    there exists a $(\psi,\mu)\in \t{Aut}(H) \times \t{Aut}(K)$ s.t. $\phi'_k = (\psi,\mu)\phi_k=\psi\circ\phi_{\mu^{-1}(k)}\circ\psi^{-1}$\\
    Let $\varphi : H \rtimes_\phi K \rightarrow H \rtimes_{\phi'} K$ s.t. $\varphi(h_1,k_1)=(\psi(h),\mu(k))$\\
    Homomorphism: Let $(h_1,k_1),(h_2,k_2) \in H \rtimes_{\phi} K$\\
    $\varphi[(h_1,k_1)(h_2,k_2)]=\varphi(h_1\phi_{k_1}(h_2),k_1k_2)=(\psi(h_1\phi_{k_1}(h_2)),\mu(k_1,k_2))$\\\
    $\varphi(h_1,k_1)\varphi(h_2,k_2)=(\psi(h_1),\mu(k_1))(\psi(h_2),\mu(k_2))=(\psi(h_1)\phi'_{\mu(k)}(\psi(h_2)),\mu(k_1)\mu(k_2))$\\
    $=(\psi(h_1)\psi\circ\phi_{\mu^{-1}(\mu(k_1))}\circ\psi^{-1}(\psi(h_2)),\mu(k_1k_2))=(\psi(h_1\phi_{k_1}(h_2)),\mu(k_1k_2))$\\
    Injective: Let $(h_1,k_1),(h_2,k_2) \in H \rtimes_{\phi} K$ s.t. $\varphi(h_1,k_1)=\varphi(h_2,k_2)$
    \begin{align*}
        \varphi(h_1,k_1)&=\varphi(h_2,k_2)\\
        (\psi(h_1),\mu(k_1))&=(\psi(h_2),\mu(k_2))
    \end{align*}
    thus $\psi(h_1)=\psi(h_2)$ and $\mu(k_1)=\mu(k_2)$ since these are automorphisms $h_1=h_2$ and $k_1=k_2$ thus $(h_1,k_1)=(h_2,k_2)$\\
    Surjective: Let $(h,k) \in H \rtimes_{\phi'} K$ consider $(\psi^{-1}(h),\mu^{-1}(k))$\\
    $\varphi(\psi^{-1}(h),\mu^{-1}(k))=(h,k)$\\\
    Thus $H \rtimes_\phi K \cong H \rtimes_{\phi'} K$
\end{enumerate}

\newpage
\subsection*{Problem 5}
Used some latex from class notes.\\
Let $n_3$ be the number of $3$-Sylows then by Sylow Theorem 3, $n_3 \mid 7$ and $n_3 \equiv 1 \pmod{3}$.
Thus, $n_3 = 1$ or $n_3 = 7$\\
Let $n_7$ be the number of $7$-Sylows then by Sylow 3, $n_7 \mid 3$ and $n_7 \equiv 1 \pmod{7}$, thus $n_7 = 1$.\\
If $n_3 = 1$ then this implies $H$ is a unique 3-Sylow subgroup of $G$, and thus it is normal by Sylow 2 this applies 
for $K$ the unique 7-Sylow subgroup of $G$.\\
Consider $HK$: Since $H,K$ are both normal $HK$ is a subgroup of $G$. Also, $H\cup K$ has $9$ elements of the 21 elements in $G$,
thus by Lagrange's theorem the only other order of elements outside $H \cup K$ would have order 21, thus $G \cong C_{21}$\\
If $n_3 = 7$. We still have that the unique $7$-Sylow subgroup $K=\{e,x,x^2,\dots,x^6\}$ which is normal in $G$, and we can fix an arbitrary $3$-Sylow subgroup $H = \{e, y, y^2 \}$.\\
Since  $K$ is normal this implies $yxy^{-1}=x^i$ thus
\begin{gather*}
    y^2 x y^{-2} = y(yxy^{-1})y^{-1} = y x^i y^{-1} = (yxy^{-1})^i = (x^i)^i = x^{(i^2)} \\
    x = y^3 x y^{-3} = y (y^2 x y^{-2})y^{-1} = y(x^{(i^2)})y^{-1} = (x^{(i^2)})^i = x^{(i^3)}
\end{gather*}
From this we can see that $i^3 \equiv 1 \pmod{7}$, reducing the possibilities to $i \in \{1, 2, 4 \}$.
\begin{itemize}
    \item If $i=1$ then $yx=xy$ thus consider $xy$
    \begin{multicols}{3}
            $(xy)^2=x^2y^2\\
            (xy)^3=x^3y^3=x^3\\
            (xy)^4=x^4y\\
            (xy)^5=x^5y^2\\
            (xy)^6=x^6\\
            (xy)^7=y\\
            (xy)^8=xy^2\\
            (xy)^9=x^2\\
            (xy)^{10}=x^3y\\
            (xy)^{11}=x^4y^2\\
            (xy)^{12}=x^5\\
            (xy)^{13}=x^6y\\
            (xy)^{14}=y^2\\
            (xy)^{15}=x\\
            (xy)^{16}=x^2y\\
            (xy)^{17}=x^3y^2\\
            (xy)^{18}=x^4\\
            (xy)^{19}=x^5y\\
            (xy)^{20}=x^6y^2\\
            (xy)^{21}=e\\
            $
    \end{multicols}
    Thus ord$(xy)=21$ thus $G \cong C_{21}$
    \item If $i=2$\\
    Thus $yxy^{-1}=x^2$\\
    Thus $G=\ang{x,y \mid x^7=y^3=e, yxy^{-1}=x^2}=A$
    \item If $i = 4$\\
    Thus $yxy^{-1}=x^4$\\
    thus $y^2xy^{-2}=x^{16}=x^2$\\
    Thus let $\varphi: G \rightarrow \ang{a,b \mid a^7=b^3=e, bab^{-1}=a^2}=A$ s.t. $\varphi(x)=a$ and $\varphi(y^2)=b$\\
    Homomorphism: We only need to check for we sent same order to same order which we did and that $\varphi(y^2)\varphi(x)=\varphi(x)^2\varphi(y^2)$\\
    $\varphi(y^2)\varphi(x)=ba=ba^2=\varphi(x)^2\varphi(y^2)$\\
    Injective: Let $g \in$ ker$(\varphi)$\\
    Since we did not send any generators of $G$ to the identity elmeent this means that our kernel is trivial thus $\varphi$ is injective\\
    Surjective: Let $g \in A$ thus $g=a^ib^j$ for $0 \leq i < 7$ and $0 \leq j < 3$\\
    Consider $x^iy^{2j}$\\
    $\varphi(x^iy^{2j})=\varphi(x)^i\varphi(y^2)^j=a^ib^j=g$\\
    Thus $G \cong A$\\
    and $A$ we know has a subgroup $H$ of size 7 that is normal, and we know it has a subgroup $K$ of size 3.\\
    Since $H$ only has elements of order 1,7 and $K$ only has elemetns of order 1,3 we know $H\cap K =\{e\}$ thus $|HK|=|H||K|/|H\cap K|=21$ thus $G=HK$ and $H \cap K = \{e\}$
    thus by 3 we know $A \cong H \rtimes K$ and since $H,K$ are prime order we know they are isomorphic to $C_7,C_3$ repectively thus $A \cong C_7 \rtimes C_3$
\end{itemize}

\newpage
\subsection*{Problem 6}
Consider: $x^n$\\
$yx^ny^{-1}=x^{2n}$ thus\\
Class 1: $\{x,x^2,x^4\}$\\
Class 2: $\{x^3,x^5,x^6\}$\\
Consider: $y$\\
$x^iyx^{-i}=yx^{2i}x^{-i}=yx^i$ for $i \in \{0,\dots,6\}$\\
Class 3: $\{y,yx,yx^2,yx^3,yx^4,yx^5,yx^6\}$\\
Consider: $y^2$\\
$x^iy^2x^{-i}=yx^{2i}yx^{-i}=y^2x^{3i}$ for $i \in \{0,\dots,6\}$\\
Class 4: $\{y^2,y^2x,y^2x^2,y^2x^3,y^2x^4,y^2x^5,y^2x^6\}$\\
Class 5: $\{e\}$\\
Thus $|G|=1+3+3+7+7=21$

\newpage
\subsection*{Problem 7}
Let $G$ be order $30=2\times3\times5$.\\
By Sylow 1 there exists 2-Sylow subgroup, 3-Sylow subgroup, and 5-Sylow subgroup.\\
Let $n_5$ be the number of 5-Sylow subgroups the by Sylow 3, $n_5 \mid 6$ and $n_5 \equiv 1 \pmod{5}$\\
Thus $n_5=1,6$\\
Let $n_3$ be the number of 3-Sylow subgroups the by Sylow 3, $n_3 \mid 10$ and $n_3 \equiv 1 \pmod{3}$\\
Thus $n_3=1,10$\\
If $n_5=6$ and $n_3=10$ then this means there are $2\cdot 10=20$ elements of order $3$\\ 
this leaves $9$ elements to be order 5 however there are $4 \cdot 6=25$ elements of order $5$ thus these cannot be true at the same time\\
Thus $n_5$ or $n_3$ is equal to 1.\\
If $n_5=1$ let $H$ be the unique 5-Sylow. Thus, $H$ is normal in $G$, and let $K$ be a 3-Sylow.\\
Since $H$ is normal this implies $HK$ is a subgroup of $G$. Also since $H$ has only elements of order 5,1 and $K$ only has elements of order 3,1 this implies $K\cap H = \{e\}$
thus $|HK|=\dfrac{|H||K|}{|H \cap K|}=15$ thus $[G : HK]=2$ by lagrange's theroem. Thus, since $HK$ is index 2 this means it is normal in $G$\\
If $n_3=1$ do the same thing as above, and you will get a normal subgroup of order 15.\\
Let $L$ be that normal subgroup of order $15$ since it is order 15 proven in class this can only be isomorphic $C_{15}$ thus let $L=\ang{x}$ for $x\in L$\\
Let $U$ be a 2-Sylow subgroup since it is order of a prime we know it is cyclic thus let $U=\ang{y}$ for $y \in U$\\
Since $L$ is normal this implies $yxy^{-1}=x^i$ thus\\
$x=y^2xy^{-2}=y(yxy^{-1})y^{-1}=yx^iy^{-1}=x^{2i}$ thus\\
$i^2 \equiv 1 \pmod{15}$\\
$i \in \{1,4,11,14\}$\\
\begin{itemize}
    \item If $i=1$ then $yx=xy$ thus consider $xy$
    \begin{multicols}{3}
        $(xy)^2=x^2y^2=x^2\\
        (xy)^3=x^3y^3=x^3y\\
        (xy)^4=x^4\\
        (xy)^5=x^5y\\
        (xy)^6=x^6\\
        (xy)^7=x^7y\\
        (xy)^8=x^8\\
        (xy)^9=x^9y\\
        (xy)^{10}=x^{10}\\
        (xy)^{11}=x^{11}y\\
        (xy)^{12}=x^{12}\\
        (xy)^{13}=x^{13}y\\
        (xy)^{14}=x^{14}\\
        (xy)^{15}=y\\
        (xy)^{16}=x\\
        (xy)^{17}=x^2y\\
        (xy)^{18}=x^3\\
        (xy)^{19}=x^4y\\
        (xy)^{20}=x^5\\
        (xy)^{21}=x^6y\\
        (xy)^{22}=x^7\\
        (xy)^{23}=x^8y\\
        (xy)^{24}=x^9\\
        (xy)^{25}=x^{10}y\\
        (xy)^{26}=x^{11}\\
        (xy)^{27}=x^{12}y\\
        (xy)^{28}=x^{13}\\
        (xy)^{29}=x^{14}y\\
        (xy)^{30}=e\\
        $
    \end{multicols}
    Thus $G \cong C_{30}$
    \item If $i=14$ then $yxy^{-1}=x^{14} \implies yx=x^{14}y$\\
    thus $G= \ang{x,y \mid x^{15}=y^2=e,\; yx=x^{14}y} \cong D_{15}$
    \item If $i=4$ then $yxy^{-1}=x^4$ $G= \ang{x,y \mid x^{15}=y^2=e,\; yxy^{-1}=x^4}$
    \item If $i=11$ then $yxy^{-1}=x^{11}$ $G= \ang{x,y \mid x^{15}=y^2=e,\; yxy^{-1}=x^{11}}$
\end{itemize}

\newpage
\subsection*{Problem 8}
Let $G$ be a group of order 12.\\
By Sylow 1 there exists 2-Sylow subgroup and a 3-Sylow subgroup.\\
Let $n_2$ be the number of 2-Sylow subgroups the by Sylow 3, $n_2 \mid 3$ and $n_2 \equiv 1 \pmod{2}$\\
Thus $n_2=1,3$\\
Let $n_3$ be the number of 3-Sylow subgroups the by Sylow 3, $n_3 \mid 4$ and $n_3 \equiv 1 \pmod{3}$\\
Thus $n_3=1,4$\\
If $n_2=3$ and $n_3=4$ then there are $2*4=8$ elements of order 3 thus there are 3 elements left for the 2-Sylow subgroups
this is not enough elements thus $n_2=1$ or $n_3=1$\\
\begin{enumerate}
    \item If $n_2=1$\\
    Then let $H$ be the 2-Sylow subgroup and $K$ be a 3-sylow subgroup and since $|K|=3$ this means $K\cong C_3$ and since $|H \cap K|=\{e\}$ since $H$ only has order 1,2,4 and $K$ only has order 1,3 so $|HK|=|H||K|=|G|$ thus $HK=G$ thus $G \cong H \rtimes K$
    \begin{enumerate}
        \item If $H \cong C_2 \times C_2$\\
        Thus $G \cong (C_2 \times C_2) \rtimes_\phi C_3$ s.t. $\phi : C_3 \rightarrow \t{Aut}(C_2 \times C_2) \cong S_3$ by homework 1 and is a homomorphism.\\
        Since all non identity elements in $C_3$ are order 3 and ord$(\phi(x)) \mid \t{ord}(x)=3$ thus $\phi(x)=id$ or equals three cycle where $x$ is the generator of $C_3$\\
        If $\phi$ is trivial then by 3c this implies that $G \cong C_2 \times C_2 \times C_3 \cong C_6 \times C_2$\\
        If $\phi$ is not trivial then we have two choices for $C_2 \times C_2 = \{e,1,2,3=12\}$ $\phi_1(x)=(123)$ and $\phi_2(x)=(132)$ since there are two 3 cycles. Consider $f : C_3 \rightarrow C_3$ s.t. $f(x)=x^{-1}$
        and $g : C_2\times C_2 \rightarrow C_2 \times C_2$ s.t. $g(x)=x$ these are both automorphisms\\
        Consider $(g,f)\phi_2(x)= g \circ \phi_2(f(x)) \circ g^{-1}=\phi_2(x^2)$\\
        $\phi_2(x^2)(1)=(132)(132)(1)=(123)(1)=2$\\
        $\phi_2(x^2)(2)=(132)(132)(2)=(123)(2)=3$\\
        Thus $(f,g)\phi_{2}(x)=\phi_1(x)$ thus by 4 since $x$ is the generator and defines the homomorphism $H \rtimes_{\phi_1} K = H \rtimes_{\phi_{2}} K$ thus $G \cong (C_2 \times C_2) \rtimes C_3$
        \item If $H=C_4$\\
        Thus $G \cong C_4 \rtimes_\phi C_3$ s.t. $\phi : C_3 \rightarrow \t{Aut}(C_4) \cong C_2$ since there are two generators of $C_4$ and is a homomorphism.\\
        non identity elements in $C_3$ are order 3 and ord$(\phi(x)) \mid \t{ord}(x)=3$ thus $\phi(x)=id$ thus by 3c $G \cong C_4 \times C_3 \cong C_{12}$ since proven in class if gcd$(a,b)=1$ then $C_a \times C_b \cong C_{ab}$
    \end{enumerate}
    \item If $n_3=1$\\
    then let $H$ be the 3-Sylow subgroup and since $|H|=3$ this means $H \cong C_3$ and let $K$ be a 2-Sylow subgroup abs since $|H \cap K|=\{e\}$ since $H$ only has order 1,3 and $K$ only has order 1,2,4 so $|HK|=|H||K|=|G|$ thus $HK=G$ thus  $G \cong H \rtimes K$
    \begin{enumerate}
        \item If $H \cong C_2 \times C_2$\\
        Then $G \cong C_3 \rtimes_{\phi} C_2 \times C_2$ s.t. $\phi : C_2 \times C_2 \rightarrow \t{Aut}(C_3) \cong C_2$ since $C_3$ has two generators and is a homomorphism.\\
        Since all non identity elements in $C_2 \times C_2=\{e,a,b,ab\}$ are order 2 and ord$(\phi(x)) \mid \t{ord}(x)=2$ thus $\phi(x)=id$ for all $x\in C_2 \times C_2$
        or $\phi_1(a)=c$ and $\phi_1(b)=c$ or $\phi_2(a)=c$ and $\phi_2(b)=id$ or $\phi_3(a)=id$ and $\phi_3(b)=c$ where $c : C_3 \rightarrow C_3$ s.t. $c(x)=x^{-1}$ for $x \in C_3$\\
        If $\phi$ is trivial then by 3c this implies that $G \cong C_3 \times C_2 \times C_2 \cong C_6 \times C_2$\\
        If $\phi$ is not trivial then consider
        \begin{itemize}
            \item $f$ swaps $a$ and $b$ and fixes $ab$
            \item $h$ swaps $a$ and $ab$ and fixes $b$
        \end{itemize}
        and $g : C_3 \rightarrow C_3$ s.t. $g(x)=x$ for all $x \in C_3$ these are all automorphisms
        \begin{itemize}
            \item Consider: $(g,f)\phi_2$
            \begin{itemize}
                \item $(g,f)\phi_2(a)=g \circ \phi_2(f(a))\circ g^{-1}=\phi_2(b)=id$
                \item $(g,f)\phi_2(b)=g \circ \phi_2(f(b))\circ g^{-1}=\phi_2(a)=c$
            \end{itemize}
            Thus $(g,f)\phi_2 = \phi_3$ thus by 4, \\$H \rtimes_{\phi_2} K \cong H \rtimes_{\phi_3} K$
            \item Consider: $(g,f)\phi_1$
            \begin{itemize}
                \item $(g,h)\phi_1(a)=g \circ \phi_1(f(a))\circ g^{-1}=\phi_1(ab)=\phi_1(a)\phi_1(b)=c \circ id = c$
                \item $(g,h)\phi_1(b)=g \circ \phi_1(f(b))\circ g^{-1}=\phi_1(b)=id$
            \end{itemize}
            Thus $(g,h)\phi_1 = \phi_2$ thus by 4, \\$H \rtimes_{\phi_1} K \cong H \rtimes_{\phi_2} K$
        \end{itemize}
        Thus $H \rtimes_{\phi_1} K \cong H \rtimes_{\phi_2} K \cong H \rtimes_{\phi_3} K$\\
        Thus $G \cong C_3 \rtimes C_2 \times C_2$
        \item If $H \cong C_4$\\
        Then $G \cong C_3 \rtimes_\phi C_4$ s.t. $\phi : C_4 \rightarrow \t{Aut}(C_3) \cong C_2$.\\
        Since the generator of $C_4$ is $4$ and ord$(\phi(x)) \mid \t{ord}(x)=4$ thus $\phi(x)=id$ or $\phi(x)=x^{-1}$\\
        If $\phi$ is trivial then by 4 $G \cong C_4 \times C_3 \cong C_{12}$\\
        If $\phi$ is not trivial then $G \cong C_3 \rtimes C_4$
    \end{enumerate}
\end{enumerate}
Thus we have $C_{12}$, $C_6 \times C_2$, $(C_2 \times C_2) \rtimes C_3$, $C_3 \rtimes C_4$, and $C_3 \rtimes (C_2 \times C_2)$\\
\begin{tabular}{c c c c c}
    \hline
    \textbf{Group} & \textbf{Abelian} & \boldmath$n_2$ & \boldmath$n_3$ & \textbf{2-Sylow} \\ \hline
    \boldmath$C_{12}$ & Yes & 1 & 1 & $C_4$ \\ \hline
    \boldmath$C_6 \times C_2$ & Yes & 1 & 1 & $C_2 \times C_2$ \\ \hline
    \boldmath$A_4$ & No & 1 & 4 & $C_2 \times C_2$ \\ \hline
    \boldmath$D_6$ & No & 3 & 1 & $C_2 \times C_2$ \\ \hline
    \boldmath$C_3 \rtimes C_4$ & No & 3 & 1 & $C_4$ \\ \hline
\end{tabular}
\phantom{text}\\
\phantom{text}\\
$(C_2 \times C_2) \rtimes C_3$ has the same properties as $A_4$ thus $(C_2 \times C_2) \rtimes C_3 \cong A_4$\\
$C_3 \rtimes (C_2 \times C_2)$ has the same properties as $D_{12}$ thus $C_3 \rtimes (C_2 \times C_2) \cong D_{12}$\\
Thus, theses $5$ groups are also not isomorphic to one another because the properties between them are different shown in the table\\
Thus, theses $5$ groups are all the groups of order 12 unique up to isomorphism.
\end{document}